% ═══════════════════════════════════════════════════════════════════════════
% section_universal_reaction_engine.tex
% Universal Automated Chemical Reaction Engine — Theory and Specification
%
% VSEPR-SIM Research Platform
% Generated: 2025
% ═══════════════════════════════════════════════════════════════════════════

\documentclass[11pt,a4paper]{article}
\usepackage[utf8]{inputenc}
\usepackage{lmodern}
\usepackage[T1]{fontenc}
\usepackage{microtype}
\usepackage{amsmath,amssymb}
\usepackage{booktabs}
\usepackage{enumitem}
\usepackage{fancyhdr}
\usepackage{geometry}
\usepackage{xcolor}
\usepackage{listings}
\usepackage{hyperref}
\usepackage[labelfont=bf]{caption}
\usepackage{titlesec}

\geometry{margin=2.5cm}
\pagestyle{fancy}
\fancyhf{}
\fancyhead[L]{\textsc{VSEPR-SIM}}
\fancyhead[R]{\textsc{Universal Reaction Engine}}
\fancyfoot[C]{\thepage}

\titleformat{\section}{\Large\bfseries}{\texorpdfstring{\S}{Sec.}\thesection}{1em}{}
\titleformat{\subsection}{\large\bfseries}{\texorpdfstring{\S}{Sec.}\thesubsection}{0.8em}{}

\definecolor{codeblue}{RGB}{40,80,160}
\lstset{
    basicstyle=\ttfamily\small,
    keywordstyle=\color{codeblue}\bfseries,
    commentstyle=\color{gray},
    frame=single,
    breaklines=true,
    numbers=left,
    numberstyle=\tiny\color{gray}
}

\title{%
    \vspace{-1cm}
    \textbf{Universal Automated Chemical Reaction Engine}\\[0.3cm]
    \large Thermal--Atomistic Simulation with Bonding Predictions\\
    and SolidWorks-Compatible Engineering Output\\[0.2cm]
    \normalsize VSEPR-SIM Research Platform
}
\author{VSEPR-SIM Development}
\date{\today}

\begin{document}
\maketitle
\thispagestyle{fancy}

% ═══════════════════════════════════════════════════════════════════════════
\section{Scope and Purpose}
\label{sec:scope}
% ═══════════════════════════════════════════════════════════════════════════

The Universal Automated Chemical Reaction Engine extends VSEPR-SIM from
atomistic structure generation into full chemical reaction simulation.  It
bridges three domains:

\begin{enumerate}[label=(\alph*)]
    \item \textbf{Chemistry} --- real stoichiometric reactions with mass balance,
          thermodynamics, and mechanism tracking.
    \item \textbf{Atomistic simulation} --- the 6+9 seed-and-bead stepper
          with Lorentz--Berthelot mixing and environment-responsive coupling.
    \item \textbf{Engineering output} --- SolidWorks-compatible CFD material
          data (12--22 element output per species), structural point clouds,
          and Excel-compatible spreadsheets.
\end{enumerate}

% ═══════════════════════════════════════════════════════════════════════════
\section{Reaction Archetypes}
\label{sec:archetypes}
% ═══════════════════════════════════════════════════════════════════════════

Three chemically distinct reaction archetypes are implemented:

\subsection{Reaction~1 --- Organic Synthesis with Inorganic Reagent}

\begin{equation}
    \text{C}_6\text{H}_6 + \text{HNO}_3
    \xrightarrow{\text{H}_2\text{SO}_4}
    \text{C}_6\text{H}_5\text{NO}_2 + \text{H}_2\text{O}
    \label{eq:nitration}
\end{equation}

Electrophilic aromatic substitution.  Benzene (organic substrate) reacts
with nitric acid (inorganic reagent) in the presence of sulfuric acid
(catalyst) to produce nitrobenzene (organic product) and water (by-product).

\subsection{Reaction~2 --- Special Inorganic Precipitation}

\begin{equation}
    \text{Th(NO}_3\text{)}_4 + 2\,\text{H}_2\text{C}_2\text{O}_4
    + 6\,\text{H}_2\text{O}
    \longrightarrow
    \text{Th(C}_2\text{O}_4\text{)}_2 \cdot 6\text{H}_2\text{O}\!\downarrow
    + 4\,\text{HNO}_3
    \label{eq:thorium}
\end{equation}

Precipitation of thorium oxalate sludge from aqueous solution.  The product
is a fine, dense precipitate used in actinide separation chemistry.

\subsection{Reaction~3 --- Thermal Decomposition (Salt \texorpdfstring{$\to$}{to} Gas + Powder)}

\begin{equation}
    2\,\text{Cu(NO}_3\text{)}_2
    \longrightarrow
    2\,\text{CuO}(s) + 4\,\text{NO}_2(g) + \text{O}_2(g)
    \label{eq:copper}
\end{equation}

Thermal decomposition of copper(II) nitrate.  Produces a black CuO powder
residue plus brown NO$_2$ gas and colourless O$_2$ gas.

% ═══════════════════════════════════════════════════════════════════════════
\section{Thermodynamic Framework}
\label{sec:thermo}
% ═══════════════════════════════════════════════════════════════════════════

All thermodynamic quantities are computed from tabulated standard enthalpies
of formation ($\Delta H_f^\circ$) and Gibbs energies ($\Delta G_f^\circ$).

\subsection{Enthalpy of Reaction}

\begin{equation}
    \Delta H_\text{rxn}^\circ
    = \sum_{\text{products}} \nu_i \, \Delta H_{f,i}^\circ
    - \sum_{\text{reactants}} \nu_j \, \Delta H_{f,j}^\circ
    \label{eq:hess}
\end{equation}

\subsection{Gibbs Energy and Equilibrium}

\begin{equation}
    \Delta G_\text{rxn}^\circ
    = \sum \nu_i \, \Delta G_{f,i}^\circ
    - \sum \nu_j \, \Delta G_{f,j}^\circ
\end{equation}

\begin{equation}
    K_\text{eq} = \exp\!\left(\frac{-\Delta G_\text{rxn}^\circ}{RT}\right)
    \label{eq:keq}
\end{equation}

\subsection{Entropy of Reaction}

\begin{equation}
    \Delta S_\text{rxn}^\circ
    = \frac{\Delta H_\text{rxn}^\circ - \Delta G_\text{rxn}^\circ}{T}
    \label{eq:entropy}
\end{equation}

\subsection{Arrhenius Kinetics}

\begin{equation}
    k(T) = A \exp\!\left(\frac{-E_a}{RT}\right)
    \label{eq:arrhenius}
\end{equation}

% ═══════════════════════════════════════════════════════════════════════════
\section{Species Data Model}
\label{sec:species}
% ═══════════════════════════════════════════════════════════════════════════

Each \texttt{ChemicalSpecies} carries:

\begin{table}[h]
\centering
\caption{Species data fields}
\label{tab:species}
\begin{tabular}{@{}lll@{}}
\toprule
Field & Units & Description \\
\midrule
\texttt{name}           & ---          & Human-readable name \\
\texttt{formula}        & ---          & Molecular formula \\
\texttt{phase}          & ---          & Solid / Liquid / Gas / Aqueous / Sludge \\
\texttt{molecular\_weight} & g/mol    & Sum of atomic masses \\
\texttt{density}        & g/cm$^3$    & Bulk density \\
\texttt{delta\_Hf}      & kcal/mol    & Standard enthalpy of formation \\
\texttt{delta\_Gf}      & kcal/mol    & Standard Gibbs energy \\
\texttt{S\_std}         & cal/(mol$\cdot$K) & Standard entropy \\
\texttt{Cp}             & cal/(mol$\cdot$K) & Heat capacity \\
\texttt{lj\_sigma}      & \AA         & LJ $\sigma$ parameter \\
\texttt{lj\_epsilon}    & kcal/mol    & LJ $\epsilon$ parameter \\
\texttt{thermal\_cond}  & W/(m$\cdot$K) & Thermal conductivity \\
\texttt{specific\_heat} & J/(g$\cdot$K) & Specific heat capacity \\
\texttt{viscosity}      & Pa$\cdot$s  & Dynamic viscosity \\
\texttt{dielectric}     & ---         & Relative permittivity \\
\texttt{dipole\_moment} & Debye       & Electric dipole moment \\
\bottomrule
\end{tabular}
\end{table}

% ═══════════════════════════════════════════════════════════════════════════
\section{Bonding Prediction}
\label{sec:bonding}
% ═══════════════════════════════════════════════════════════════════════════

Bonding predictions combine:

\begin{enumerate}[label=\arabic*.]
    \item \textbf{Electronegativity difference} --- Pauling's rule:
          $\Delta\text{EN} > 1.7 \Rightarrow$ ionic character.
    \item \textbf{Bond change tracking} --- explicit record of every bond
          broken and formed, with energies.
    \item \textbf{Environment-responsive state} --- $\eta$ at the time of
          bond formation/breaking provides a structural context measure.
\end{enumerate}

% ═══════════════════════════════════════════════════════════════════════════
\section{Mass Balance}
\label{sec:balance}
% ═══════════════════════════════════════════════════════════════════════════

Every reaction is verified for mass conservation:

\begin{equation}
    \sum_{\text{reactants}} \nu_j \cdot n_{j,\text{elem}}
    = \sum_{\text{products}} \nu_i \cdot n_{i,\text{elem}}
    \quad \forall \; \text{elements}
\end{equation}

Tolerance: $|\Delta| < 10^{-6}$ moles per element.  By-products are
tracked explicitly and included in the balance.

% ═══════════════════════════════════════════════════════════════════════════
\section{SolidWorks Engineering Output}
\label{sec:solidworks}
% ═══════════════════════════════════════════════════════════════════════════

\subsection{CFD Material Data (22-Element CSV)}

Each species produces a 22-field record suitable for SolidWorks Flow
Simulation import:

\begin{table}[h]
\centering
\caption{22-element material property output}
\label{tab:cfd}
\begin{tabular}{@{}clll@{}}
\toprule
\# & Field & Units & Category \\
\midrule
1  & Species name      & ---          & Identity \\
2  & Formula           & ---          & Identity \\
3  & Phase             & ---          & Identity \\
4  & Density           & g/cm$^3$    & Mechanical \\
5  & Molecular weight  & g/mol       & Mechanical \\
6  & Viscosity         & Pa$\cdot$s  & Mechanical \\
7  & Surface tension   & N/m         & Mechanical \\
8  & Temperature       & K           & Thermal \\
9  & Thermal cond.     & W/(m$\cdot$K) & Thermal \\
10 & Specific heat     & J/(g$\cdot$K) & Thermal \\
11 & Cp                & J/(mol$\cdot$K) & Thermal \\
12 & Melting point     & K           & Thermal \\
13 & Boiling point     & K           & Thermal \\
14 & $\Delta H_f$      & kcal/mol    & Thermodynamic \\
15 & $\Delta G_f$      & kcal/mol    & Thermodynamic \\
16 & $S^\circ$         & cal/(mol$\cdot$K) & Thermodynamic \\
17 & $\Delta H_\text{rxn}$ & kcal/mol & Thermodynamic \\
18 & $\varepsilon_r$   & ---         & Electrical \\
19 & $\mu$ (dipole)    & Debye       & Electrical \\
20 & $\sigma_\text{LJ}$ & \AA        & Atomistic \\
21 & $\varepsilon_\text{LJ}$ & kcal/mol & Atomistic \\
22 & $\langle\text{EN}\rangle$ & Pauling & Atomistic \\
\bottomrule
\end{tabular}
\end{table}

\subsection{Material XML (.sldmat)}

An XML material library file compatible with SolidWorks material import:

\begin{lstlisting}[language=XML,caption={SolidWorks material XML format}]
<material name="Benzene">
  <physicalproperties>
    <DENS value="876.5" units="kg/m^3"/>
    <KX value="0.159" units="W/(m*K)"/>
    <C value="1740" units="J/(kg*K)"/>
    <MU value="0.000604" units="Pa*s"/>
  </physicalproperties>
</material>
\end{lstlisting}

% ═══════════════════════════════════════════════════════════════════════════
\section{Implementation Summary}
\label{sec:implementation}
% ═══════════════════════════════════════════════════════════════════════════

\begin{table}[h]
\centering
\caption{Source file inventory}
\label{tab:files}
\begin{tabular}{@{}lll@{}}
\toprule
File & Purpose & Lines \\
\midrule
\texttt{chemistry/species.hpp}         & Species data model      & $\sim$210 \\
\texttt{chemistry/reaction.hpp}        & Reaction framework      & $\sim$260 \\
\texttt{chemistry/reaction\_library.hpp} & 3 built-in reactions  & $\sim$490 \\
\texttt{chemistry/reaction\_engine.hpp}  & Simulation + export   & $\sim$440 \\
\texttt{example\_universal\_reaction\_engine.cpp} & Main demo    & $\sim$190 \\
\texttt{section\_universal\_reaction\_engine.tex} & This document & $\sim$280 \\
\bottomrule
\end{tabular}
\end{table}

% ═══════════════════════════════════════════════════════════════════════════
\section{Dependencies}
\label{sec:deps}
% ═══════════════════════════════════════════════════════════════════════════

\begin{table}[h]
\centering
\caption{Module dependencies}
\label{tab:deps}
\begin{tabular}{@{}ll@{}}
\toprule
Module & Depends On \\
\midrule
\texttt{species.hpp}          & \texttt{atomistic/core/state.hpp} \\
\texttt{reaction.hpp}         & \texttt{species.hpp} \\
\texttt{reaction\_library.hpp} & \texttt{reaction.hpp} \\
\texttt{reaction\_engine.hpp}  & \texttt{reaction\_library.hpp}, \texttt{seed\_bead\_stepper.hpp}, \\
                                & \texttt{snapshot\_graph.hpp}, \texttt{excel\_export.hpp}, \\
                                & \texttt{solidworks\_export.hpp} \\
\bottomrule
\end{tabular}
\end{table}

\end{document}
